% Exploration: twelve ways to draw the feasible solid F (a thickened band) and its carving from C.
\documentclass[10pt,border=2mm]{standalone}
\usepackage[T1]{fontenc}
\usepackage{figurestyle}
\input{primitives.tex}
\input{numbers.tex}
\newcommand\cell[3]{\node[sub,anchor=north,text width=3.5cm,align=center] at (#1,#2) {#3};}
\begin{document}
\begin{tikzpicture}[plate]
\path[use as bounding box] (0,0) rectangle (15.6,16.8);
\pgfmathsetmacro{\Ro}{1.1}\pgfmathsetmacro{\Ri}{.78}\pgfmathsetmacro{\Rm}{.94}\pgfmathsetmacro{\Hh}{.3}\pgfmathsetmacro{\E}{.42}
% ---------- 1 wedge cut (current)
\begin{scope}[shift={(2.0,15.0)}]\tikzset{tube/Ro=1.1,tube/Ri=.78,tube/H=.3,tube/E=.42,tube/a1=-62,tube/a2=8}\pic (a) at (0,0) {thicktube};\foreach \p in {top0,top120,top240} {\hatomat{a-\p}{.06}{ColGold}}\end{scope}
\cell{2.0}{13.6}{1. wedge cut, band at mid-thickness (current)}
% ---------- 2 quarter removed
\begin{scope}[shift={(5.9,15.0)}]\tikzset{tube/Ro=1.1,tube/Ri=.78,tube/H=.3,tube/E=.42,tube/a1=-90,tube/a2=0}\pic (b) at (0,0) {thicktube};\foreach \p in {top120,top240} {\hatomat{b-\p}{.06}{ColGold}}\end{scope}
\cell{5.9}{13.6}{2. a quarter removed: more of the inside, one version lost to the cut}
% ---------- 3 sliced loaf: eight transverse slices with gaps
\begin{scope}[shift={(9.8,15.0)}]
  \foreach \k in {0,...,7} {\pgfmathsetmacro{\ta}{-90+45*\k+6}\pgfmathsetmacro{\tb}{-90+45*\k+39}
    \path[fill=ColNitrogen!16] plot[domain=\ta:\tb,samples=10,variable=\a] ({\Ro*sin(\a)},{\Hh-\E*\Ro*cos(\a)}) -- plot[domain=\tb:\ta,samples=10,variable=\a] ({\Ro*sin(\a)},{-\Hh-\E*\Ro*cos(\a)}) -- cycle;
    \foreach \a in {\ta,\tb} {\path[fill=white] ({\Ri*sin(\a)},{\Hh-\E*\Ri*cos(\a)})--({\Ro*sin(\a)},{\Hh-\E*\Ro*cos(\a)})--({\Ro*sin(\a)},{-\Hh-\E*\Ro*cos(\a)})--({\Ri*sin(\a)},{-\Hh-\E*\Ri*cos(\a)})--cycle;
      \path[cut] ({\Ri*sin(\a)},{\Hh-\E*\Ri*cos(\a)})--({\Ro*sin(\a)},{\Hh-\E*\Ro*cos(\a)})--({\Ro*sin(\a)},{-\Hh-\E*\Ro*cos(\a)})--({\Ri*sin(\a)},{-\Hh-\E*\Ri*cos(\a)})--cycle;
      \draw[fine] ({\Ri*sin(\a)},{\Hh-\E*\Ri*cos(\a)})--({\Ro*sin(\a)},{\Hh-\E*\Ro*cos(\a)})--({\Ro*sin(\a)},{-\Hh-\E*\Ro*cos(\a)})--({\Ri*sin(\a)},{-\Hh-\E*\Ri*cos(\a)})--cycle;
      \draw[ColNitrogen,line width=.8pt] ({\Rm*sin(\a)},{\Hh-\E*\Rm*cos(\a)})--({\Rm*sin(\a)},{-\Hh-\E*\Rm*cos(\a)});}
    \path[fill=white] plot[domain=\ta:\tb,samples=10,variable=\a] ({\Ro*sin(\a)},{\Hh-\E*\Ro*cos(\a)}) -- plot[domain=\tb:\ta,samples=10,variable=\a] ({\Ri*sin(\a)},{\Hh-\E*\Ri*cos(\a)}) -- cycle;
    \draw[fine] plot[domain=\ta:\tb,samples=10,variable=\a] ({\Ro*sin(\a)},{\Hh-\E*\Ro*cos(\a)}) -- plot[domain=\tb:\ta,samples=10,variable=\a] ({\Ri*sin(\a)},{\Hh-\E*\Ri*cos(\a)}) -- cycle;
    \draw[fine] plot[domain=\ta:\tb,samples=10,variable=\a] ({\Ro*sin(\a)},{-\Hh-\E*\Ro*cos(\a)});}
\end{scope}
\cell{9.8}{13.6}{3. a sliced loaf: F as a circle of normal slices, every cut face shows the band}
% ---------- 4 ribbon of cross-sections on the band
\begin{scope}[shift={(13.7,15.0)}]\tikzset{band/R=.94,band/H=.3,band/E=.42}\pic (d) at (0,0) {band};
  \foreach \a in {-60,0,60} {\path[cut] ({\Ri*sin(\a)},{\Hh-\E*\Ri*cos(\a)})--({\Ro*sin(\a)},{\Hh-\E*\Ro*cos(\a)})--({\Ro*sin(\a)},{-\Hh-\E*\Ro*cos(\a)})--({\Ri*sin(\a)},{-\Hh-\E*\Ri*cos(\a)})--cycle;
    \draw[fine] ({\Ri*sin(\a)},{\Hh-\E*\Ri*cos(\a)})--({\Ro*sin(\a)},{\Hh-\E*\Ro*cos(\a)})--({\Ro*sin(\a)},{-\Hh-\E*\Ro*cos(\a)})--({\Ri*sin(\a)},{-\Hh-\E*\Ri*cos(\a)})--cycle;}
\end{scope}
\cell{13.7}{13.6}{4. the band with a few normal slices standing across it: thickness as a family of slices}
% ---------- 5 wireframe tube around the band
\begin{scope}[shift={(2.0,10.9)}]\tikzset{band/R=.94,band/H=.3,band/E=.42}\pic (e) at (0,0) {band};
  \draw[fine] (0,\Hh) ellipse[x radius=\Ro,y radius={\E*\Ro}]; \draw[fine] (0,\Hh) ellipse[x radius=\Ri,y radius={\E*\Ri}];
  \draw[fine] plot[domain=-90:90,samples=30,variable=\a] ({\Ro*sin(\a)},{-\Hh-\E*\Ro*cos(\a)}); \draw[hidden] plot[domain=90:270,samples=30,variable=\a] ({\Ro*sin(\a)},{-\Hh-\E*\Ro*cos(\a)});
  \draw[fine] (-\Ro,-\Hh)--(-\Ro,\Hh) (\Ro,-\Hh)--(\Ro,\Hh);
  \foreach \a in {-60,-30,0,30,60} \draw[line width=.25pt,draw=PlateInk] ({\Ro*sin(\a)},{\Hh-\E*\Ro*cos(\a)})--({\Ro*sin(\a)},{-\Hh-\E*\Ro*cos(\a)});
\end{scope}
\cell{2.0}{9.5}{5. the band solid, the wall as a wireframe around it: the neighbourhood as an outline only}
% ---------- 6 plan view
\begin{scope}[shift={(5.9,10.9)}]
  \path[fill=ColNitrogen!16] (0,0) circle[radius=\Ro]; \path[fill=white] (0,0) circle[radius=\Ri];
  \path[cut] (0,0) -- (-62:\Ro) arc[start angle=-62,end angle=8,radius=\Ro] -- cycle; \path[fill=white] (0,0) circle[radius=\Ri];
  \draw[heavy] (0,0) circle[radius=\Ro]; \draw[heavy] (0,0) circle[radius=\Ri]; \draw[hidden] (0,0) circle[radius=\Rm];
  \draw[heavy] (-62:\Ri)--(-62:\Ro) (8:\Ri)--(8:\Ro);
  \draw[tpath] plot[domain=90:-30,samples=20,variable=\a] ({\Rm*cos(\a)},{\Rm*sin(\a)});
  \foreach \a in {90,-30,210} {\hatom{\Rm*cos(\a)}{\Rm*sin(\a)}{.06}{ColGold}}
\end{scope}
\cell{5.9}{9.5}{6. plan view, 2D on purpose: annulus, band circle, versions, the removed sector hatched}
% ---------- 7 exploded carving: the slab with the hole, the solid lifted out
\begin{scope}[shift={(9.8,10.9)}]
  \path[cut] (-1.9,-1.1)--(1.3,-1.1)--(1.9,-.5)--(-1.3,-.5)--cycle; \draw[heavy] (-1.9,-1.1)--(1.3,-1.1)--(1.9,-.5)--(-1.3,-.5)--cycle;
  \path[fill=white,even odd rule] (0,-.8) ellipse[x radius=\Ro,y radius={\E*\Ro}] (0,-.8) ellipse[x radius=\Ri,y radius={\E*\Ri}];
  \draw[fine] (0,-.8) ellipse[x radius=\Ro,y radius={\E*\Ro}]; \draw[fine] (0,-.8) ellipse[x radius=\Ri,y radius={\E*\Ri}];
  \tikzset{tube/Ro=1.1,tube/Ri=.78,tube/H=.3,tube/E=.42,tube/a1=-62,tube/a2=8}\pic (g) at (0,.9) {thicktube};
  \draw[hidden] (-\Ro,-.8)--(-\Ro,.6) (\Ro,-.8)--(\Ro,.6);
\end{scope}
\cell{9.8}{9.5}{7. exploded carving: the solid lifted out of its hole in the hatched slab $\conf$}
% ---------- 8 window in the wall
\begin{scope}[shift={(13.7,10.9)}]\tikzset{tube/Ro=1.1,tube/Ri=.78,tube/H=.3,tube/E=.42,tube/a1=0,tube/a2=0}\pic (h) at (0,0) {thicktube};
  \begin{scope}\clip ({\Ro*sin(-40)},{-.16-\E*\Ro*cos(-40)})--({\Ro*sin(-8)},{-.16-\E*\Ro*cos(-8)})--({\Ro*sin(-8)},{.16-\E*\Ro*cos(-8)})--({\Ro*sin(-40)},{.16-\E*\Ro*cos(-40)})--cycle;
    \path[fill=white] (-2,-2) rectangle (2,2);
    \foreach \a in {94,100,...,266} \draw[line width=.25pt,draw=PlateInk!70] ({\Ri*sin(\a)},{\Hh-\E*\Ri*cos(\a)})--({\Ri*sin(\a)},{-\Hh-\E*\Ri*cos(\a)});
    \draw[ColNitrogen,line width=.8pt] plot[domain=-45:-3,samples=12,variable=\a] ({\Rm*sin(\a)},{-\E*\Rm*cos(\a)});\end{scope}
  \draw[heavy] ({\Ro*sin(-40)},{-.16-\E*\Ro*cos(-40)})--({\Ro*sin(-8)},{-.16-\E*\Ro*cos(-8)})--({\Ro*sin(-8)},{.16-\E*\Ro*cos(-8)})--({\Ro*sin(-40)},{.16-\E*\Ro*cos(-40)})--cycle;
\end{scope}
\cell{13.7}{9.5}{8. a window cut in the wall: the band seen inside; the solid stays whole}
% ---------- 9 half tube: split along a vertical plane, halves drawn apart
\begin{scope}[shift={(2.0,6.8)}]
  \foreach \dx/\ta/\tb in {-.35/-180/0,.35/0/180} {
    \begin{scope}[shift={(\dx,0)}]
      \path[fill=ColNitrogen!16] plot[domain=\ta:\tb,samples=30,variable=\a] ({\Ro*sin(\a)},{\Hh-\E*\Ro*cos(\a)}) -- plot[domain=\tb:\ta,samples=30,variable=\a] ({\Ro*sin(\a)},{-\Hh-\E*\Ro*cos(\a)}) -- cycle;
      \path[fill=white] plot[domain=\ta:\tb,samples=30,variable=\a] ({\Ro*sin(\a)},{\Hh-\E*\Ro*cos(\a)}) -- plot[domain=\tb:\ta,samples=30,variable=\a] ({\Ri*sin(\a)},{\Hh-\E*\Ri*cos(\a)}) -- cycle;
      \draw[fine] plot[domain=\ta:\tb,samples=30,variable=\a] ({\Ro*sin(\a)},{\Hh-\E*\Ro*cos(\a)}) -- plot[domain=\tb:\ta,samples=30,variable=\a] ({\Ri*sin(\a)},{\Hh-\E*\Ri*cos(\a)}) -- cycle;
      \foreach \a in {\ta,\tb} {\path[cut] ({\Ri*sin(\a)},{\Hh-\E*\Ri*cos(\a)}) rectangle ({\Ro*sin(\a)},{-\Hh-\E*\Ro*cos(\a)}); \draw[fine] ({\Ri*sin(\a)},{\Hh-\E*\Ri*cos(\a)}) rectangle ({\Ro*sin(\a)},{-\Hh-\E*\Ro*cos(\a)}); \draw[ColNitrogen,line width=.8pt] ({\Rm*sin(\a)},{\Hh-\E*\Rm*cos(\a)})--({\Rm*sin(\a)},{-\Hh-\E*\Rm*cos(\a)});}
    \end{scope}}
\end{scope}
\cell{2.0}{5.4}{9. split along the axis plane, halves drawn apart: four cut faces, the band on each}
% ---------- 10 machine drawing: hidden lines, no tint, dense generatrix shading, section hatching
\begin{scope}[shift={(5.9,6.8)}]
  \foreach \a in {-88,-84,...,88} {\pgfmathsetmacro{\lw}{0.15+0.45*abs(sin(\a))^3}\draw[line width=\lw pt,draw=PlateInk] ({\Ro*sin(\a)},{\Hh-\E*\Ro*cos(\a)})--({\Ro*sin(\a)},{-\Hh-\E*\Ro*cos(\a)});}
  \draw[hidden] plot[domain=90:270,samples=30,variable=\a] ({\Ro*sin(\a)},{-\Hh-\E*\Ro*cos(\a)}); \draw[hidden] (0,-\Hh) ellipse[x radius=\Ri,y radius={\E*\Ri}];
  \draw[heavy] plot[domain=-90:90,samples=30,variable=\a] ({\Ro*sin(\a)},{-\Hh-\E*\Ro*cos(\a)}); \draw[heavy] (-\Ro,-\Hh)--(-\Ro,\Hh) (\Ro,-\Hh)--(\Ro,\Hh);
  \path[fill=white] (0,\Hh) ellipse[x radius=\Ro,y radius={\E*\Ro}]; \draw[heavy] (0,\Hh) ellipse[x radius=\Ro,y radius={\E*\Ro}];
  \path[fill=white] (0,\Hh) ellipse[x radius=\Ri,y radius={\E*\Ri}]; \draw[heavy] (0,\Hh) ellipse[x radius=\Ri,y radius={\E*\Ri}];
  \foreach \a in {94,100,...,266} {\pgfmathsetmacro{\lw}{0.15+0.4*abs(cos(\a))}\draw[line width=\lw pt,draw=PlateInk] ({\Ri*sin(\a)},{\Hh-\E*\Ri*cos(\a)})--({\Ri*sin(\a)},{\Hh-\E*\Ri*cos(\a)-.0001});}
  \draw[hidden] (0,\Hh) ellipse[x radius=\Rm,y radius={\E*\Rm}];
\end{scope}
\cell{5.9}{5.4}{10. machine drawing: no tint, hidden lines dashed, generatrix shading only (the hole floor hidden)}
% ---------- 11 orthographic triple: plan, elevation, section
\begin{scope}[shift={(9.8,6.8)}]
  \begin{scope}[shift={(-1.0,.35)}]\draw[fine] (0,0) circle[radius=.55]; \draw[fine] (0,0) circle[radius=.39]; \draw[hidden] (0,0) circle[radius=.47]; \path[cut] (0,0)--(-62:.55) arc[start angle=-62,end angle=8,radius=.55]--cycle; \path[fill=white] (0,0) circle[radius=.39]; \draw[fine] (-62:.39)--(-62:.55) (8:.39)--(8:.55);\end{scope}
  \begin{scope}[shift={(.6,.35)}]\draw[fine] (-.55,-.15) rectangle (.55,.15); \draw[hidden] (-.55,0)--(.55,0); \path[cut] (-.55,-.15) rectangle (-.39,.15); \path[cut] (.39,-.15) rectangle (.55,.15); \draw[fine] (-.39,-.15)--(-.39,.15) (.39,-.15)--(.39,.15); \draw[ColNitrogen,line width=.8pt] (-.47,-.15)--(-.47,.15) (.47,-.15)--(.47,.15);\end{scope}
  \begin{scope}[shift={(-.2,-.75)}]\draw[fine] (-.55,-.15) rectangle (.55,.15); \draw[hidden] (-.39,-.15)--(-.39,.15) (.39,-.15)--(.39,.15); \draw[hidden] (-.47,-.15)--(-.47,.15) (.47,-.15)--(.47,.15);\end{scope}
  \node[sub] at (-1.0,-.4) {plan}; \node[sub] at (.6,-.4) {section}; \node[sub] at (-.2,-1.1) {elevation};
\end{scope}
\cell{9.8}{5.4}{11. orthographic triple: plan, section through the axis, elevation; engineering convention}
% ---------- 12 layer cake: the eta-foliation, five thin annular discs
\begin{scope}[shift={(13.7,6.8)}]
  \foreach \k in {-2,...,2} {\pgfmathsetmacro{\y}{.17*\k}
    \path[fill=ColNitrogen!16] plot[domain=-90:90,samples=30,variable=\a] ({\Ro*sin(\a)},{\y+.05-\E*\Ro*cos(\a)}) -- plot[domain=90:-90,samples=30,variable=\a] ({\Ro*sin(\a)},{\y-.05-\E*\Ro*cos(\a)}) -- cycle;
    \draw[fine] plot[domain=-90:90,samples=30,variable=\a] ({\Ro*sin(\a)},{\y-.05-\E*\Ro*cos(\a)});
    \path[fill=white,even odd rule] (0,\y+.05) ellipse[x radius=\Ro,y radius={\E*\Ro}] (0,\y+.05) ellipse[x radius=\Ri,y radius={\E*\Ri}];
    \draw[fine] (0,\y+.05) ellipse[x radius=\Ro,y radius={\E*\Ro}]; \draw[fine] (0,\y+.05) ellipse[x radius=\Ri,y radius={\E*\Ri}];}
\end{scope}
\cell{13.7}{5.4}{12. layer cake: the solid as a stack of thin annuli (slices of constant $\eta$)}
\node[sub,anchor=north west,text width=15cm] at (.3,3.6) {What each teaches (notes after rendering, docs/illustration-notes.md). The band must stay visible as the object the versions live on; the thickness must read as a neighbourhood, not a pipe; the cut must read as carving.};
\end{tikzpicture}
\end{document}
